\chapter{User and Group Administration \& Privilege Architecture}
\label{chap:user-group-admin}

\section{The Linux Multi-User Administrative Hierarchy}
User identity in Linux forms the security perimeter. The kernel does not identify entities by human-readable usernames; internally, all authorization is mapped to numeric 32-bit \textbf{User Identifiers (UIDs)} and \textbf{Group Identifiers (GIDs)}.

\begin{definition}[Classification of User Accounts]
\begin{enumerate}[leftmargin=1.5em]
    \item \textbf{Superuser (root):} Assigned $\text{UID} = 0$. The root account bypasses all filesystem permission checks, can bind to privileged TCP/UDP network ports ($<1024$), manipulate raw disk sectors, and signal any process.
    \item \textbf{System Service Accounts:} Typically assigned UIDs between $1$ and $999$ (e.g., \texttt{bin}, \texttt{daemon}, \texttt{www-data}, \texttt{sshd}). These accounts own background service daemons and have non-interactive login shells (\texttt{/usr/sbin/nologin} or \texttt{/bin/false}).
    \item \textbf{Regular Human Users:} Assigned UIDs starting at $1000$ (e.g., \texttt{1000}, \texttt{1001}). Represent actual human developers and interactive operators.
\end{enumerate}
\end{definition}

\subsection{Primary versus Supplementary Groups}
Every user belongs to exactly one \textbf{Primary Group} (recorded in their \texttt{/etc/passwd} entry). Files newly created by the user automatically inherit this primary group. A user may additionally belong to multiple \textbf{Supplementary (Secondary) Groups} (recorded in \texttt{/etc/group}) to obtain access to shared team folders, administrative utilities (\texttt{sudo}), or hardware peripherals (\texttt{docker}, \texttt{vboxusers}).

\section{Core Account Databases: /etc/passwd and /etc/shadow}
User account properties are stored within colon-delimited ASCII configuration files.

\subsection{Structure of /etc/passwd}
World-readable (\texttt{chmod 644}) because utilities such as \texttt{ls -l} require access to translate numeric UIDs to usernames.
\begin{center}
\texttt{username : password : UID : GID : GECOS : home\_directory : login\_shell}
\end{center}
\begin{enumerate}[leftmargin=1.5em]
    \item \texttt{username}: Unique account login string.
    \item \texttt{password}: Stored as an \texttt{x}, indicating that the actual cryptographic hash is isolated inside \texttt{/etc/shadow}.
    \item \texttt{UID}: Numeric user identifier.
    \item \texttt{GID}: Numeric primary group identifier.
    \item \texttt{GECOS / Comment}: Metadata (full user name, office location, contact phone).
    \item \texttt{home\_directory}: Initial directory upon interactive login.
    \item \texttt{login\_shell}: Absolute path to the user's login shell program.
\end{enumerate}

\subsection{Structure of /etc/shadow}
Restricted to root access exclusively (\texttt{chmod 640} or \texttt{600}) to prevent offline dictionary and rainbow table attacks.
\begin{center}
\texttt{username : hash : last\_change : min\_days : max\_days : warn\_days : inact : expire : flag}
\end{center}
\begin{itemize}[leftmargin=1.5em]
    \item \texttt{hash}: Cryptographic hash. Commences with an algorithm identifier (e.g., \texttt{\$6\$} indicates SHA-512 with a salt string).
    \item \texttt{last\_change}: Days elapsed since Epoch (Jan 1, 1970) that the password was last modified.
    \item \texttt{min\_days}: Minimum number of days required before the user can change password again.
    \item \texttt{max\_days}: Password expiration interval (forcing periodic rotation).
    \item \texttt{warn\_days}: Advance warning days provided to the user prior to password expiration.
\end{itemize}

\subsection{Structure of /etc/group}
\begin{center}
\texttt{group\_name : group\_password : GID : user\_list}
\end{center}

\section{Account Provisioning Commands}
\begin{commandbox}{useradd --- Create a New User Account}
\textbf{Syntax:} \texttt{useradd [options] username}
\begin{itemize}[leftmargin=1.5em]
    \item \texttt{-m}: Create the user's home directory (copying default profiles from \texttt{/etc/skel}).
    \item \texttt{-d /path}: Custom home directory path.
    \item \texttt{-s /path/shell}: Set default login shell (e.g., \texttt{/bin/bash}).
    \item \texttt{-g group}: Explicit primary group assignment.
    \item \texttt{-G group1,group2}: Comma-separated list of supplementary groups.
    \item \texttt{-c "Full Name"}: GECOS comment field.
    \item \texttt{-e YYYY-MM-DD}: Account expiration date.
\end{itemize}
\end{commandbox}

\begin{commandbox}{usermod --- Modify User Account Properties}
\textbf{Syntax:} \texttt{usermod [options] username}
\begin{itemize}[leftmargin=1.5em]
    \item \texttt{-aG group}: \textbf{Append} user to a supplementary group.
    \begin{warningbox}[The Fatal usermod -G Trap]
    Always pair \texttt{-G} with \texttt{-a} (\textbf{append})! If \texttt{usermod -G docker user} is run without \texttt{-a}, the user is \textbf{removed from all other existing supplementary groups}, including \texttt{sudo}, instantly revoking their administrative privileges!
    \end{warningbox}
    \item \texttt{-s /path}: Modify login shell.
    \item \texttt{-L}: Lock user account (prepends \texttt{!} to password hash in \texttt{/etc/shadow}).
    \item \texttt{-U}: Unlock a locked user account.
\end{itemize}
\end{commandbox}

\begin{commandbox}{userdel --- Delete User Account}
\textbf{Syntax:} \texttt{userdel [options] username}
\begin{itemize}[leftmargin=1.5em]
    \item \texttt{userdel username}: Removes account from \texttt{/etc/passwd} and \texttt{/etc/shadow}, leaving home directory untouched.
    \item \texttt{userdel -r username}: Recursively removes the user's home directory and mail spool.
\end{itemize}
\end{commandbox}

\section{Group Management Commands}
\begin{itemize}[leftmargin=1.5em]
    \item \texttt{groupadd -g GID groupname}: Allocates a new group with an optional explicit numeric GID.
    \item \texttt{groupmod -n newname oldname}: Renames an existing group.
    \item \texttt{groupdel groupname}: Removes an existing group (prohibited if designated as a primary group for any user).
\end{itemize}

\section{Password Administration \& Aging Policies}
\begin{itemize}[leftmargin=1.5em]
    \item \texttt{passwd [username]}: Sets or updates account passwords.
    \item \texttt{passwd -l username}: Locks user password.
    \item \texttt{passwd -e username}: Immediately expires password, forcing rotation upon next login.
    \item \textbf{chage (Change Age):} Configures enterprise password aging policies:
    \begin{itemize}
        \item \texttt{chage -l username}: Lists active password expiration and aging policy details.
        \item \texttt{chage -M 90 -m 7 -W 14 username}: Sets maximum lifespan to 90 days, minimum interval to 7 days, with 14 days warning.
    \end{itemize}
\end{itemize}

\section{Privilege Escalation: sudo and visudo}
Direct login as root is considered an operational anti-pattern. Contemporary Linux infrastructures delegate administrative privileges via \textbf{sudo} (SuperUser DO).
\begin{itemize}[leftmargin=1.5em]
    \item \textbf{Configuration in /etc/sudoers:} Edited exclusively via \textbf{visudo}, which performs atomic syntax locking and validation, preventing syntax corruption that could lock all administrators out of the operating system.
    \item \textbf{Grammar of Sudoers Rules:}
    \begin{center}
    \texttt{user\_or\_group \quad host=(runas\_user:runas\_group) \quad commands}
    \end{center}
    \textit{Example rule:} \texttt{\%devops ALL=(ALL:ALL) ALL} grants all members of the \texttt{devops} group full root escalation rights on all hosts.
    \item \textbf{su vs su -:}
    \begin{itemize}
        \item \texttt{su}: Switches user but preserves the environment variables and working directory of the original caller.
        \item \texttt{su -} (Login Shell): Completely clears environment, loading the target user's profile, \texttt{PATH}, and home directory.
    \end{itemize}
\end{itemize}

\section{Chapter Review \& Analytical Questions}
\begin{enumerate}[leftmargin=1.5em]
    \item Explain why the encrypted password field in \texttt{/etc/passwd} contains an \texttt{x} instead of the actual cryptographic hash.
    \item Contrast the impact of executing \texttt{usermod -G sudo rahul} versus \texttt{usermod -aG sudo rahul}.
    \item Decode the meaning of the 7 colon-delimited fields for the account \texttt{git:x:998:998:git daemon:/var/lib/git:/usr/bin/git-shell}.
    \item What is the functional difference between locking an account via \texttt{passwd -l} and setting its shell to \texttt{/sbin/nologin}?
    \item Why must administrators use \texttt{visudo} rather than standard text editors to update \texttt{/etc/sudoers}?
\end{enumerate}
